\subsection{饱和子集}

设 \(\mathbb{R}\{x, y\}\) 是集合 \(\mathbf{E}\) 上的一个等价关系，并且设 \(\mathbf{A}\) 为 \(\mathbf{E}\) 的一个子集。称 A 关于 \(\mathbb{R}\) 是饱和的，如果关系 \(x \in \mathbf{A}\) （关于 \(x\) ) 与 \(\mathbb{R}\{x, y\}\) 相容；或者等价地，如果对每个 \(x \in \mathbf{A}\) ， \(x\) 的等价类包含于 \(\mathbf{A}\) 。换句话说，一个集合关于 \(\mathbb{R}\) 是饱和的当且仅当它是关于 \(\mathbb{R}\) 的一组等价类的并集.

设 \(f\) 是从 \(\mathbf{E}\) 到上 \(\mathbf{E} / \mathbf{R}\) 的典范映射。如果 \(\mathbf{A}\) 关于 \(\mathbf{R}\) 是饱和的，则每个 \(x \in \mathbf{A}\) 的等价类（它等于 \(\overline{f}^1 \langle \{f(x)\} \rangle\) ）包含于 \(\mathbf{A}\) ，因此 \(\overline{f}^1 \langle f\langle \mathbf{A}\rangle \rangle \subset \mathbf{A}\) ；因为无论如何我们都有 \(\mathbf{A} \subset \overline{f}^1 \langle f\langle \mathbf{A}\rangle \rangle\) ，由此可得 \(\overline{f}^1 \langle f\langle \mathbf{A}\rangle \rangle = \mathbf{A}\) 。反之，若 \(\mathbf{A} = \overline{f}^1 \langle f\langle \mathbf{A}\rangle \rangle\) ，那么对于每一个 \(x \in \mathbf{A}\) ，等价类 \(\mathbf{K} = f(x)\) （ \(x\) 关于 \(\mathbf{R}\) 的等价类）是 \(f\langle \mathbf{A}\rangle\) 的一个元素；并且由于 \(\mathbf{K} = \overline{f}^1 \langle \{\mathbf{K}\}\rangle\) ，我们有 \(\mathbf{K} \subset \overline{f}^1 \langle f\langle \mathbf{A}\rangle \rangle = \mathbf{A}\) 。因此， \(\mathbf{E}\) 的关于 \(\mathbf{R}\) 饱和的子集恰好就是那些子集 \(\mathbf{A}\) （属于 \(\mathbf{E}\) ，使得 \(\mathbf{A} = \overline{f}^1 \langle f\langle \mathbf{A}\rangle\rangle\) ）。等价地，它们是 \(\mathbf{E}\) 的形如 \(\overline{f}^1\langle\mathbf{B}\rangle\) （其中 \(\mathbf{B} \subset E / R\) ）的子集；因为关系 \(A = \overline{f}^1\langle B\rangle\) 蕴含 \(B = f\langle A\rangle\) ，因此 \(A = \overline{f}^1\langle f\langle A\rangle\rangle\) .

如果 \((\mathbf{X}_t)_{t\in I}\) 是 \(E\) 的一族饱和子集，则集合 \(\bigcup_{t\in I} X_t\) 并且 \(\bigcap_{t\in I} X_t\) 是饱和的（§4，命题3和4）。如果 \(A\) 是 \(E\) ，那么也有 \(\complement_E A\) 的饱和子集（§4，命题6）。

现在设 A 是 E 的任意子集。则集合 \(\overline{f}^1\langle f\langle \mathrm{A}\rangle \rangle\) 包含 A 且是饱和的。反之，如果 \(\mathbf{A}'\) 是包含 A 的 E 的饱和子集，则有 \(f\langle \mathrm{A}'\rangle \supset f\langle \mathrm{A}\rangle\) ，因此

\[
\mathrm{A} ^ {\prime} = \overline {{f}} ^ {1} \langle f \langle \mathrm{A} ^ {\prime} \rangle \rangle \supset \overline {{f}} ^ {1} \langle f \langle \mathrm{A} \rangle \rangle .
\]

因此我们可以说 \(\overline{f}^1\langle f\langle A\rangle \rangle\) 是包含 A 的 E 的“最小”饱和子集（参见第三章）；这个集合称为 A 关于关系 R 的饱和化。容易看出，A 的饱和化是 A 中元素的等价类的并集。如果 \((\mathbf{X}_{i})_{i\in I}\) 是 E 的一族子集，且 \(A_{i}\) 是 \(X_{i}\) 关于 R 的饱和化，则 \(\bigcup_{i\in I}X_{i}\) 是 \(\bigcup_{i\in I}A_{i}\) 的饱和化（§4，命题3）。

